\section{Detection independent of the current width}
\label{ap:relative-detection}

We prove the warm clause of Theorem~\ref{ap:path:detection}.
The layer is $j$ and the testing degree is $r<j$. The desired
detection and approximation bounds and accuracies are warm;
auxiliary Fourier comparisons and transfer errors may use the
stated cold and late costs. The coefficient law and its late
parameters are the same as in Section~\ref{ap:paths}; no new
conditioning on productive paths is made.

Lemma~\ref{ap:path:comparison} writes the sampled cube average as
a bounded combination of ambient cube averages whose site
factors are bounded by a warm multiple of $m_j^{-1}$. Applying an inverse
theorem directly to those factors would charge the reciprocal
chart mass to the output. We instead use high rank to control
products of the chart weight, replace all but one of the sparse
factors by bounded functions, and apply the inverse theorem to
bounded conditional averages. The resulting detector and
approximation have costs independent of $Q_j$.

\subsection{A pseudorandom chart weight}

Let $\mathcal H=\prod_a\Z/(4H_a\Z)$ and use representative
coordinates to evaluate $C$ outside $U$.  An additive cube all of
whose vertices lie in $U$ has no wrap in its two-face identities;
the conversion between box and group normalizations costs only
dimension factors.  Use the cutoff $\chi$ chosen in
Section~\ref{ap:paths}, equal to one on a neighborhood of all
twists and actual path supports. Define
\begin{equation}\label{ap:path:weight}
 a_\chi=\int_{W/\Lambda}\chi,\qquad
 \nu(u)=\frac{\chi(C(u)-c)}{a_\chi}.
\end{equation}
Here the cutoff is read through the buffered ambient chart.  By
Lemma~\ref{ap:path:chart}, $a_\chi\asymp_{\rm warm}m_j$.
Every sparse factor in Equation~\eqref{ap:path:sparse-comparison}
has absolute value at most $B\nu$ for warm $B$.

\begin{lemma}[Bounded-grid moments]\label{ap:path:grid}
For $1\le q\le j$ and any bounded Cartesian grid of distinct choice
vectors $S$, the late rank and side choices can ensure
\begin{equation}\label{ap:path:grid-moment}
 \E\prod_{\alpha\in S}
 \nu\left(\sum_{a=1}^q x_{a,\alpha_a}\right)=1+o(1),
\end{equation}
where all $x_{a,b}$ are independent uniform variables of
$\mathcal H$.  The error can be made smaller than any specified
pre bound, simultaneously for the bounded grid enlargements used
below.  In particular $\E\nu=1+o(1)$.  Marginal spatial boundary
and smooth chart estimates obey the same permitted cold costs.
\end{lemma}

\begin{proof}
Fourier approximate the normalized cutoffs at cold cost in their
ambient charts.  A nontrivial resulting character has phase
\[
 \sum_{\alpha\in S}\lambda_\alpha
       C\left(\sum_a x_{a,\alpha_a}\right),
\]
where at least one row is nonzero on $W$; rows annihilating $W$
may be deleted.  Its degree-$j$ tensors are independent in the
required sense.  Indeed fix $\alpha$, substitute
$x_{a,\alpha_a}=t_a z$ and all other $x_{a,b}=0$, and extract
the coefficient of $t_1^{j-q+1}t_2\cdots t_q$.  All other
choice vectors disappear, and the result is
\[
 \frac{j!}{(j-q+1)!}\lambda_\alpha C_{[j]}(z).
\]
Thus a slow-plus-rational certificate for the joint symbol gives
one for a nonzero row of $C_{[j]}$, at cold cost.

To apply this to the cyclic representatives, partition wrap cells
into sufficiently fine boxes and discard their boundary strips.
All caps, Fourier masses, and inverse cell sizes are charged at
cold scale.  On each cell the wrap changes are constant
translations and leave the preceding homogeneous calculation
unchanged.  If a nontrivial character had nonnegligible bias,
Lemma~\ref{prelim:step-drop} on a last-layer torus, followed by
the coefficient extraction just given and restoration of the box
scales, would contradict rank $R_j$.  The trivial character gives
the product of Haar means, namely one.  Choosing the errors before
the late cutoffs proves the estimate.  Inserting a prescribed
smooth spatial weight, or approximating a boundary strip, gives
the marginal assertions by the same argument.  Actual coincidences
of randomly sampled sites are included in these averages; distinctness
here concerns the formal choice vectors.
\end{proof}

\subsection{Densification without paying the chart mass}

We use a cube version of the densification method of
Conlon, Fox and Zhao~\cite[Sections~6.2--6.3]{ConlonFoxZhao2015}.
The bounded-grid estimates above supply the required estimates
for products of the chart weight. We retain the details because
the conclusion concerns normalized chart measure and warm
niltest complexity.

\begin{lemma}[Relative cube detection]\label{ap:path:relative-cube}
Suppose a mixed $q$-cube average as in
Equation~\eqref{ap:path:sparse-comparison} has absolute value at
least $\delta$, and write its factors as
$F_\omega=P_\omega J_\omega/m_j$, where $|P_\omega|\le1$
and $J_\omega$ are warm normalized site twists supported where
$\chi=1$.  Suppose $|F_\omega|\le B\nu$, where $1\le q\le j$
is fixed and $B,\delta^{-1}$ are warm-bounded.  Fix a detector
vertex $s$.  Its input $P_s$ correlates, in normalized chart measure, with
a warm normalized twist times an ordinary degree-$(q-1)$ niltest,
by a warm inverse amount.  For $q=1$ the niltest can be constant.
\end{lemma}

\begin{proof}
Extend all factors by zero from $U$ to $\mathcal H$. The group
cube average is the box cube average multiplied by
\[
 \kappa=\frac{|\operatorname{Cub}_q(U)|}{|\mathcal H|^{q+1}}.
\]
Both $\kappa$ and its reciprocal have warm dimension bounds.
Absorb this factor into the correlation threshold and work with
group cubes throughout the proof.

Conjugations can be absorbed into the factors.  First replace all
vertices other than the designated detector by unit-bounded
functions, one at a time.  At a vertex $i$, write the average as
$\E F_i D_i$, where $D_i$ is the conditional average of the
other factors given that site.  Every current factor has cap
$B'(1+\nu)$, with warm $B'$.  Cauchy--Schwarz gives
\[
 \E(1+\nu)|D_i|^2
 \ge \frac{\delta^2}{(B')^2\E(1+\nu)}.
\]
We claim
\begin{equation}\label{ap:path:densification-moments}
 \E(\nu-1)|D_i|^2=o(1),\qquad
 \E|D_i|^4=O_{\rm warm}(1).
\end{equation}
For the fourth moment expand four cubes with a common root.
Represent them by sums of $q$ inputs, sharing the root input on
each axis and using a separate alternate input for each cube.
Every nonroot vertex has a copy-specific alternate coordinate,
so these vertices have distinct choice vectors.  Their product
majorant is bounded by Equation~\eqref{ap:path:grid-moment}.

For the first estimate use two such rooted cubes.  Every nonroot
factor misses at least one of the $q$ root inputs.  Apply weighted
Cauchy--Schwarz successively to these root inputs.  At a step,
put each factor independent of the duplicated input outside the
inner average under its majorant $B'(1+\nu)$.  A previously
introduced positive majorant is carried once at this step, not
squared.  A dependent factor is copied to distinct choices of
the duplicated input.  Distinctness persists: descendants of
different original vertices remain distinct after forgetting
the clone labels, and descendants of one vertex differ in a
cloned coordinate on which that vertex depends.  Every nonroot
descendant retains an alternate coordinate, so never coincides
with a root descendant.  After $q$ steps the copies of $\nu-1$
occupy distinct root-only choices and all other factors have
been majorized on distinct nonroot choices.  The outer masses
are bounded by Equation~\eqref{ap:path:grid-moment}.  Expanding
the remaining $(\nu-1)$ and $(1+\nu)$ factors, the main terms
cancel, while their errors are arbitrarily small.  Fixed-depth
Cauchy--Schwarz proves the first estimate in
Equation~\eqref{ap:path:densification-moments}.

Consequently $\E|D_i|^2$ is warm-large.  Interpolation yields
\[
 \E|D_i|\ge
 \frac{(\E|D_i|^2)^{3/2}}{(\E|D_i|^4)^{1/2}}.
\]
Choosing a unit phase in place of $F_i$ preserves a warm-large
cube average.  At this stage no pointwise cap for a conditional
dual containing several sparse factors has been asserted.

Now only the detector vertex $s$ is sparse.  For another vertex
$i\ne s$, conditional on its site $u$, the sparse site is
\[
 u+\sum_{a=1}^q(s_a-i_a)h_a.
\]
One coefficient is $1$ or $-1$, so this site is exactly uniform
on $\mathcal H$.  Dropping the other bounded factors proves
the pointwise estimate
\begin{equation}\label{ap:path:dual-cap}
 |D_i(u)|\le B\E_{\mathcal H}\nu=O_{\rm warm}(1).
\end{equation}
This uses pairwise uniformity of group cube sites, not a
pointwise conditional pseudorandomness assertion for $\nu$.
Zero extensions, localization, and representative wraps are
included in the bounded factors and do not change the estimate.

If the current correlation is at least $\delta'$, then
$\E|D_i|^2\ge(\delta')^2$.  The sparse factor satisfies
\[
 \|F_s\|_{U^q}^{2^q}
 \le B^{2^q}\E_{\rm cube}\prod_\omega\nu(u_\omega)
 =O_{\rm warm}(1)
\]
by Lemma~\ref{ap:path:grid}.  Inserting the appropriate conjugate
of $D_i$ at vertex $i$ and applying Gowers--Cauchy--Schwarz gives
\[
 (\delta')^2\le \|D_i\|_{U^q}\|F_s\|_{U^q}
                    \prod_{\omega\ne i,s}\|F_\omega\|_{U^q}.
\]
Thus $D_i$ has a warm-large Gowers norm. For $q>1$, apply the form of
Theorem~\ref{prelim:inverse} for products of cyclic groups on
$\mathcal H$ to the bounded function in
Equation~\eqref{ap:path:dual-cap}, after warm rescaling.  It
gives a correlating ordinary degree-$(q-1)$ niltest.  Replacing
the factor at $i$ by that niltest preserves a warm correlation;
repeat at each nondetector vertex.  The inverse theorem has
never been applied to a function with cap $m_j^{-1}$.

The representative-coordinate box version of this step is
obtained by partitioning coordinates into short representative
intervals, using the Gowers triangle inequality to choose a
large piece, and smoothing its boundaries, as in
Theorem~\ref{prelim:inverse}.  The dual inputs are bounded.
After all replacements, fix the increments to obtain a
correlation of the sparse detector with a product of translated
niltests.  Wrap and localization boundaries in this last
selection can be smoothed using the spatial and chart marginal
bounds of Lemma~\ref{ap:path:grid}.  Products and linear
localizing masks have the permitted degree and warm complexity.
When $q=1$, the two group sites are independent, so a large
cube average already gives a large detector mean and the same
conclusion with the constant niltest.

Finally the sparse twist is supported where the cutoff in
Equation~\eqref{ap:path:weight} is one.  Dividing that factor
by $\nu$ there changes it only by a warm normalization.  The
detector is supported in $U$. Restricting the group average to
$U$ and normalizing by $\E_U\nu$ therefore gives the measure
$\mu_c$ defined in Section~\ref{ap:paths}; these normalizations
have warm cost. The resulting twist is unit-bounded after warm
rescaling, as claimed.
\end{proof}

\begin{lemma}[Weighted detection implies approximation]
\label{ap:path:weighted-separation}
Lemma~\ref{cube:separation} remains valid with any positive
probability measure $\mu$ on a finite set in place of its
uniform measure.  Its coefficient-mass and atom-count bounds
depend only on the same quantitative cap, tail, and detection
data, not on the smallest atom of $\mu$.
\end{lemma}

\begin{proof}
Remove zero-mass points.  Suppose
$X(g)\le C\E_\mu m|g|$, $\E_\mu m\le C$, and the
cap and detector hypotheses hold.  Write
$B=CC_0/\eta$ and $\delta_0=\delta(B,\eta/2)$, and request
\[
 \E_\mu m1_{m>C_0}
 \le\frac\eta C\min(1/8,\delta_0/4).
\]
For $M>\max(1,2/\delta_0)$, if $(9/8)f$ is outside
$M\overline{\operatorname{aconv}}(\mathcal J)$ plus the
$X$-ball of radius $\eta$, finite-dimensional separation in
the $\mu$ pairing gives $\chi$ with
\[
 \Re\langle(9/8)f,\chi\rangle_\mu=1,\quad
 X^*(\chi)\le\eta^{-1},\quad
 \sup_{J\in\mathcal J}|\langle J,\chi\rangle_\mu|
 \le M^{-1}.
\]
Testing at a single point gives $|\chi|\le Cm/\eta$;
the positive atom mass cancels from this inequality.  Radially
clip $\chi$ at $B$ to obtain $v$.  The tail bound gives
$\E_\mu|\chi-v|\le\min(1/8,\delta_0/4)$ and
$\E_\mu|v|\ge55/72$.  Hence
\[
 X(v)\ge\eta\Re\langle v,\chi\rangle_\mu
 \ge\eta\E_\mu|v|^2>\eta/2.
\]
Detection contradicts the last bound on the pairings with
$\chi$.  This proves the convex-hull approximation with
$X$-error at most $8\eta/9$.  Approximate the closed hull
finitely, then sample its unit-bounded atoms with their phases.
The expected squared $L^2((1+m)\mu)$ sampling error is at
most $(1+C)M^2/N$, and
$X(g)\le C^{3/2}\|g\|_{L^2((1+m)\mu)}$.
Choosing $N$ from $C,M,\eta$ uses the remaining error margin.
This proves the assertion without dependence on atom masses.
\end{proof}

\begin{proof}[Proof of the warm clause of Theorem~\ref{ap:path:detection}]
Mesh and enumerate the allowed pathwise slices as in the cold
proof. Lemma~\ref{prelim:dual-gowers} and
Lemma~\ref{ap:path:comparison} give ambient cube detection at
order $q=r+1\le j$. Lemma~\ref{ap:path:relative-cube} therefore
gives warm detection in normalized chart measure.

Here is the precise change of measure for approximation.  Let
$m_U$ be the full one-site density relative to uniform measure
on $U$, and set
\[
 d\mu=d\mu_c=\frac{\nu}{\E_U\nu}\,dU,
 \qquad m_{\rm rel}=(\E_U\nu)\frac{m_U}{\nu}.
\]
Sites of zero $\nu$ mass are ignored; no path or test support
uses them.  On actual supports $\nu\asymp_{\rm warm}m_j^{-1}$,
and $\E_U\nu=1+o(1)$ by rank.  The cap and excess-mass
statement in Lemma~\ref{ap:path:ensemble} therefore gives a
warm cap and arbitrarily small excess mass for $m_{\rm rel}$ under
$\mu$.  Enlarge that cap from $C_0$ to $2C_0$; the inequality
$m_{\rm rel}1_{m_{\rm rel}>2C_0}
 \le2(m_{\rm rel}-C_0)_+$ gives the required arbitrarily small
tail.  If every permitted slice has at least fraction
$\delta_{\rm sl}$ of full parameter mass, then
\[
 X_r(g)\le\delta_{\rm sl}^{-1}\E_U m_U|g|
       =\delta_{\rm sl}^{-1}\E_\mu m_{\rm rel}|g|.
\]
This bound holds before taking the supremum; it exchanges no
supremum with conditional expectation.  The detector of
Lemma~\ref{ap:path:relative-cube} is a $\mu$ pairing with an
allowed twist times niltest.  Apply
Lemma~\ref{ap:path:weighted-separation}, choosing its cap and
detection request first and the one necessary tail accuracy
afterwards.  All are warm requests, whereas the late choices
make the required tail as small as needed.  This proves the
full arbitrary-input approximation assertion and its stated
uniformities.
\end{proof}


